% Shared preamble for every EXAM template (midterm, final, theory, question bank,
% cheat sheet). Edit templates/_exam_preamble.tex, then run: python3 scripts/sync_preamble.py
% That stamps it between the ">>> SHARED PREAMBLE" markers of each exam template.
% \documentclass, geometry and the "who and what" block stay per template.
\usepackage{setspace}
\usepackage{amsmath,amssymb}
\usepackage{xcolor}
\usepackage{titlesec}
\usepackage{enumitem}
\usepackage{booktabs}
\usepackage{array}
\usepackage{multicol}
\usepackage{needspace}
\usepackage{silence}
\WarningFilter{LaTeX}{Command \showhyphens has changed}
\usepackage{microtype}
\usepackage{fancyhdr}
% Fonts: tectonic runs XeTeX, so use fontspec (OpenType Latin Modern covers
% accented names like "Trần Bảo Thành"). pdflatex falls back to the classic setup.
\usepackage{iftex}
\ifTUTeX
  \usepackage{fontspec}
\else
  \usepackage[T1]{fontenc}
  \usepackage[utf8]{inputenc}
  \usepackage{lmodern}
\fi
\usepackage{parskip}
\usepackage{listings}
\usepackage{tikz}
\usetikzlibrary{arrows.meta,positioning,calc,fit,shapes.geometric,backgrounds}
% bookmarks=false: math in a heading + PDF bookmarks hangs the compile forever
\usepackage[hidelinks,bookmarks=false]{hyperref}
\usepackage{tcolorbox}
\tcbuselibrary{skins,breakable}

% ---- spacing and page breaks ----
\setstretch{1.05}
\setlength{\parskip}{5pt}
\setlength{\parindent}{0pt}
\setlist[itemize]{itemsep=2pt,topsep=3pt,parsep=0pt}
\setlist[enumerate]{itemsep=2pt,topsep=3pt,parsep=0pt}
\raggedbottom
\widowpenalty=10000 \clubpenalty=10000
\displaywidowpenalty=10000 \predisplaypenalty=10000
\setlength{\emergencystretch}{2em}

% ---- section formatting (a heading never sits alone at the bottom of a page) ----
\titleformat{\section}{\large\bfseries}{}{0pt}{\needspace{6\baselineskip}}[\titlerule]
\titleformat{\subsection}{\normalsize\bfseries}{}{0pt}{\needspace{5\baselineskip}}
\titleformat{\subsubsection}{\normalsize\itshape\bfseries}{}{0pt}{\needspace{4\baselineskip}}
\titlespacing{\section}{0pt}{14pt}{6pt}
\titlespacing{\subsection}{0pt}{10pt}{4pt}
\titlespacing{\subsubsection}{0pt}{6pt}{2pt}

% ---- part banner: \examPart{Title}{subtitle} ----
\newcommand{\examPart}[2]{\clearpage\begin{center}{\Large\bfseries #1}\\[3pt]{\normalsize #2}\end{center}\vspace{2pt}\hrule\vspace{6pt}}

% ---- math shortcuts (add your own) ----
\newcommand{\R}{\mathbb R}
\newcommand{\sign}{\operatorname{sign}}
\newcommand{\diag}{\operatorname{diag}}

% ---- small inline helpers ----
% \sref{L1 p46}            inline source tag, lighter than a box title
% \gl{term:dot}            superscript link back to a glossary item with \label{term:dot}
% \qtag{LIKELY}            tag badge: LIKELY, PAST EXAM A, PRACTICE, COMPUTE, TWIST, ...
% \exam{Midterm + Final}   which exam a section is for (put it in the heading)
% \lbl{What it is:}        bold run-in label for fixed-pattern blocks
% \examcue{"quote"}        what the lecturer stressed, verbatim
\newcommand{\sref}[1]{\textcolor{blue!50!black}{\small[#1]}}
\newcommand{\gl}[1]{\textsuperscript{\textcolor{blue!50!black}{[\ref{#1}]}}}
\newcommand{\qtag}[1]{{\setlength{\fboxsep}{1.5pt}\colorbox{black!10}{\scriptsize\bfseries\MakeUppercase{#1}}}}
\newcommand{\exam}[1]{\ \textnormal{\small(#1)}}
\newcommand{\lbl}[1]{\par\noindent\textbf{#1}\ }
\newcommand{\examcue}[1]{\par\noindent\textcolor{red!55!black}{\textbf{Exam cue:}}\ \emph{#1}\par}

% ---- plain boxes: white background, coloured streak on the left ----
\tcbset{
  streak/.style={
    enhanced, breakable, colback=white, frame hidden, boxrule=0pt, arc=0pt,
    left=8pt, right=3pt, top=3pt, bottom=3pt,
    borderline west={3pt}{0pt}{#1},
    before skip=8pt, after skip=8pt
  },
  streaktitle/.style 2 args={
    streak=#1, title={#2}, attach title to upper, after title={\par\smallskip},
    fonttitle=\small\bfseries\color{#1}
  },
  titled/.style 2 args={
    enhanced, colback=#1!2, colframe=#1!45!black, boxrule=0.6pt, arc=2pt,
    left=7pt, right=7pt, top=5pt, bottom=5pt, before skip=10pt, after skip=10pt,
    colbacktitle=#1!45!black, coltitle=white, fonttitle=\bfseries, title={#2}
  },
  bw/.style={
    enhanced, breakable, colback=white, colframe=black, boxrule=0.4pt, arc=2pt,
    left=4pt, right=4pt, top=2pt, bottom=2pt, before skip=6pt, after skip=6pt
  }
}
\newtcolorbox{qbox}{streak=black!70}
\newtcolorbox{resultbox}{streak=black!45}
\newtcolorbox{trapbox}{streaktitle={orange!80!black}{Traps}}
\newtcolorbox{algobox}[1]{streaktitle={violet!70!black}{#1}}
% defbox: the lecturer's exact wording for a definition / named list. \begin{defbox}{Lecture 2, slide 7}
\newtcolorbox{defbox}[1]{streaktitle={teal!60!black}{Lecturer's definition \textnormal{\small(#1)}}}
% proofbox: a proof written the way it should appear on the paper. \begin{proofbox}{Title}
\newtcolorbox{proofbox}[1]{titled={orange}{Proof: #1\hfill{\normalfont\small copy this structure}}, breakable}

% ---- question box: \begin{origbox}{Title}{Source, e.g. Pool \S2.1, p.~1} ----
% Not breakable on purpose: a question should never split across pages.
\newtcolorbox{origbox}[2]{titled={blue}{#1\hfill{\normalfont\small #2}}}

% ---- guideline box: the recipe, written with letters so it works for any numbers ----
\newtcolorbox{guidebox}{titled={red}{Step-by-step guideline\hfill{\normalfont\small works for any numbers}}, breakable}

% ---- answer box: \begin{examsolbox}{Question 2.1} ----
\newtcolorbox{examsolbox}[1]{titled={green}{Answer: #1\hfill{\normalfont\small what to write on the exam paper}}, breakable}

% ---- theory-course boxes ----
% notebox     black frame, print friendly (overview, decoders, case studies)
% memobox     "Memorise this" ordered list / chain, with an optional title
% scenariobox a situational / essay question with its answer skeleton
% summarybox  the last-minute checklist at the end of every sheet
% pitfallbox  exam-day pitfalls, gathered in one numbered list
\newtcolorbox{notebox}{bw}
\newtcolorbox{memobox}[1][Memorise this]{streaktitle={teal!60!black}{#1}}
\newtcolorbox{scenariobox}[1]{titled={violet}{#1\hfill{\normalfont\small situational question}}, breakable}
\newtcolorbox{summarybox}{titled={black}{Quick-reference summary\hfill{\normalfont\small last-minute checklist}}, breakable}
\newtcolorbox{pitfallbox}{titled={orange}{Exam-day pitfalls\hfill{\normalfont\small read before you start}}, breakable}

% ---- question bank: \Q{Q12}{\qtag{past exam A} Birthday attack on SHA-2} ----
% then \begin{qbankbox} question only \end{qbankbox}, then Answer / Why / Trap / Also asked as
\newcommand{\Q}[2]{\needspace{6\baselineskip}\subsection*{#1\quad #2}}
\newtcolorbox{qbankbox}{streak=blue!50!black}

% ---- step and iteration badges ----
% \step{1}{Lagrangian}   -> red "STEP 1" badge + bold name on its own line
% \iterhead{2}{at $(0.8,0)$:} -> blue "ITERATION 2" badge
\newcommand{\step}[2]{\par\medskip\noindent{\setlength{\fboxsep}{2pt}\colorbox{red!55!black}{\textcolor{white}{\small\bfseries\,STEP #1\,}}}\ \ {\bfseries #2}\par\nobreak\smallskip}
\newcommand{\iterhead}[2]{\par\medskip\noindent{\setlength{\fboxsep}{2pt}\colorbox{blue!45!black}{\textcolor{white}{\small\bfseries\,ITERATION #1\,}}}\ \ #2\par\nobreak\smallskip}

% ---- code listings (black and white, prints well) ----
\lstdefinestyle{pyStyle}{language=Python, basicstyle=\ttfamily\small, keywordstyle=\bfseries,
  commentstyle=\itshape, stringstyle=\ttfamily, breaklines=true, frame=single, numbers=left,
  numberstyle=\tiny, xleftmargin=15pt, captionpos=b, columns=fullflexible, keepspaces=true}
\lstdefinestyle{cppStyle}{style=pyStyle, language=C++}

% ---- TikZ styles for redrawing exam diagrams (block ciphers, pipelines, flows) ----
\tikzset{
  blk/.style={draw, minimum height=6.5mm, minimum width=11mm, font=\scriptsize, align=center, inner sep=2pt},
  enc/.style={draw, thick, minimum height=9mm, minimum width=18mm, font=\scriptsize\bfseries, align=center},
  xr/.style={draw, circle, minimum size=4.5mm, inner sep=0pt, font=\footnotesize},
  kb/.style={draw, fill=black!10, minimum height=7mm, minimum width=15mm, font=\scriptsize, align=center, inner sep=2pt},
  lb/.style={font=\scriptsize, align=center},
  ar/.style={-{Stealth[length=1.6mm]}}
}

\pagestyle{fancy}
\fancyhf{}
\cfoot{\small\thepage}
\renewcommand{\headrulewidth}{0pt}
